% Canonical template for every problem record in database/problems/.
%
% Rules:
% 1. Copy this structure into database/problems/<op_id>.tex, where <op_id> is
%    the stable ID printed by `node scripts/new-problem-id.mjs`.
% 2. Put only the descriptive title in \section; the site never shows a
%    problem number, so do not prefix the title with "Problem N".
% 3. Put every displayed expression in a numbered equation environment, give
%    it a unique label of the form eq:<prefix>-description, and cite it with
%    \eqref. Choose a short prefix unique to this record (for example the last
%    four characters of the ID).
% 4. Use record-specific reference labels of the form ref:<prefix>-description.
% 5. Make each problem statement independent and self-contained. Describe
%    relationships to other problems only in the Comment subsection.
% 6. Use alpha-style author--year labels for references and citations.
% 7. Place Source between Status and Progress. Cite the paper or preprint that
%    states the problem explicitly, or identify concise implicit sources. If no
%    literature source exists, give the contributor's name; use "unknown" only
%    when the contributor is also unknown.
% 8. Keep Tag and ID as the final two subsections, in that order. Assign one
%    to six canonical tags from database/tags.json, separated by semicolons.
% 9. Run `node site/build.mjs` before committing; the build rejects records
%    that break these rules.

\section{Short descriptive title}
\label{sec:problem-PREFIX}

\paragraph{Problem.}
State the problem here without relying on another problem's definitions or
notation.  Introduce each displayed equation in the prose, then use a numbered
and labeled equation environment.  Refer back to every display using
\eqref{eq:PREFIX-description}.
\begin{equation}
  f(x)=x,
  \label{eq:PREFIX-description}
\end{equation}

\subsection*{Status}

% Use exactly one of: Solved, Unsolved, Partially solved.
Unsolved

\subsection*{Source}

Briefly identify the paper or preprint that posed the problem, and state
whether the attribution is explicit or implicit
\sourcecite{ref:PREFIX-source}{AutYY}.  Multiple sources may be cited when
needed.  If no literature source exists, write \texttt{Contributor: Full
Name.}; if the contributor is unknown, write \texttt{unknown}.

\subsection*{Progress}

\begin{itemize}
  \item Summarize a known result and cite its source with
  \sourcecite{ref:PREFIX-source}{AutYY}.
\end{itemize}

\subsection*{References}

\begin{enumerate}[label={},leftmargin=4.8em,labelsep=0.5em]
  \footnotesize
  \item[\textup{[AutYY]}]\label{ref:PREFIX-source}
  Author, ``Title,'' \emph{Journal} \textbf{volume}, pages (year).
  \href{https://doi.org/example}{doi:example};
  \href{https://arxiv.org/abs/example}{arXiv:example}.
\end{enumerate}

\subsection*{Comment}

State precisely what remains unresolved.  If this problem is related to
another problem in the zoo, explain that relationship here rather than in the
problem statement.

\subsection*{Tag}

% Use one to six names from database/tags.json exactly as written and
% separate multiple tags with semicolons.
Canonical tag one; Canonical tag two

\subsection*{ID}

% Paste the output of `node scripts/new-problem-id.mjs` and escape the
% underscore for LaTeX.  Never reuse or regenerate an assigned ID.
\texttt{op\_0123456789abcdef}
